\begin{lemma}{Sortierungsmengen}{Sortierungsmengen}
Seien $n \in \N$ und $x \in \R^n$.\\
Für jedes $r \in \haken n$ sei
$ M_{r,x}:=\menge{x_i}{i\in\haken n,\abs{\menge{j\in\haken n}{x_j>x_i}}<r}$.\\
Dann gilt:
\begin{enumerate}[(1)]
	\item \label{Sortierungsmengen1}
	$\forall\, k \in \haken {n-1}: M_{k, x} \subseteq M_{k+1,x}$.
	\item \label{Sortierungsmengen2}
	$\forall\, k \in \haken n: M_{k, x} \neq \emptyset$.
\end{enumerate}
\end{lemma}
\begin{proof}
(1) Sei $k \in \haken n$ und $y \in M_{k, x}$.\\
Dann gibt es ein $i \in \haken n$ mit $y = x_i$ und
$\abs{\menge{j\in\haken n}{x_j>x_i}}<k < k+1$.\\
Also ist $y=x_i \in M_{k+1,x}$.\\
(2) Wegen (1) genügt es zu zeigen, dass $M_{1, x} \neq \emptyset$ gilt.\\
Es gilt:
\begin{align*}
M_{1,x} &= \menge{x_i}{i\in\haken n,\abs{\menge{j\in\haken n}{x_j>x_i}}<1}
= \menge{x_i}{i\in\haken n, \menge{j\in\haken n}{x_j>x_i} = \emptyset}\\
&= \menge{x_i}{i\in\haken n, \forall\, j \in \haken n: x_i \ge x_j}
= \menge{x_i}{i \in \haken n, x_i = \max\menge{x_j}{j\in \haken n}}\text.
\end{align*}
Da $\menge{x_j}{j \in \haken n}$ endlich und nichtleer ist, gibt es
$l \in \haken n$ mit $x_l = \max\menge{x_j}{j \in \haken n}$ und damit ist
$M_{1, x}$ wegen obiger Umformung nichtleer.
\end{proof}